\section{Introducción}

El embarazo ectópico es un tipo de embarazo patológico que se da cuando el embrión se implanta fuera de la cavidad uterina, pudiendo ocasionar complicaciones graves y siendo una de las causas principales de morbimortalidad materna en todo el mundo. En aproximadamente el 95\% de los casos, la implantación tiene lugar en las trompas de Falopio, lo que supone un riesgo de rotura tubárica y hemorragias graves \cite{Mullany2023, Papageorgiou2025}.

Entre los principales factores de riesgo se encuentran los antecedentes de embarazo ectópico, la enfermedad inflamatoria pélvica, las intervenciones quirúrgicas abdominales y pélvicas, la ligadura de trompas previa al embarazo, y la infertilidad. Otros factores de riesgo con una menor asociación son la edad materna avanzada, el tabaquismo, la existencia de múltiples parejas sexuales, y los antecedentes de abortos espontáneos e inducidos \cite{Brim2025}. Sin embargo, también se presentan casos en pacientes sin factores de riesgo conocidos, por lo que resulta de interés estudiar los mecanismos biológicos que pueden intervenir en su aparición \cite{Mullany2023}.

En condiciones fisiológicas, tras la fecundación, el embrión es transportado a través de las trompas de Falopio hasta la cavidad intrauterina, donde tiene lugar la implantación. En este proceso intervienen las contracciones de la musculatura lisa, el movimiento del fluido tubárico y la actividad de los cilios móviles presentes en el epitelio de las trompas. 

Los cilios son estructuras celulares cuyo movimiento coordinado contribuye al transporte del contenido tubárico. Las alteraciones en su formación o funcionamiento pueden afectar a la función reproductiva y dificultar el transporte embrionario. De hecho, estudios recientes han observado una disminución de las células multiciliadas en mujeres con embarazo ectópico, lo que sugiere una posible relación entre  las alteraciones ciliares y esta patología \cite{Jiang2022}.

Asimismo, estudios experimentales en modelos animales han mostrado que determinadas alteraciones genéticas que afectan a la motilidad ciliar pueden provocar retrasos en el transporte de los ovocitos \cite{Scully2025}. Estos hallazgos ponen de manifiesto el interés de estudiar los genes implicados en la función ciliar y su posible relación con el embarazo ectópico.
